\documentclass[11pt,letterpaper]{article}
\usepackage[margin=0.88in,headheight=16pt]{geometry}
\usepackage{fontspec}
\setmainfont{Latin Modern Roman}
\setsansfont{DejaVu Sans}
\setmonofont{DejaVu Sans Mono}[Scale=0.84]
\usepackage{amsmath,amssymb,mathtools}
\usepackage{booktabs,longtable,array,tabularx}
\usepackage[table]{xcolor}
\usepackage{enumitem,fancyhdr,fvextra}
\usepackage{xurl}
\usepackage{hyperref}
\definecolor{ink}{HTML}{19354A}
\definecolor{accent}{HTML}{176D83}
\definecolor{pale}{HTML}{EFF5F7}
\hypersetup{colorlinks=true,linkcolor=accent,urlcolor=accent,
  pdftitle={Shortening and simplifying the DiophantineConsistency Lean proofs},
  pdfauthor={Research and refactoring review prepared for Vladimir Reshetnikov}}
\pagestyle{fancy}
\fancyhf{}
\fancyhead[L]{\small\sffamily ProveIt / Lean proof simplification}
\fancyhead[R]{\small\sffamily 26 October 2026}
\fancyfoot[C]{\small\thepage}
\renewcommand{\headrulewidth}{0.3pt}
\setlength{\parindent}{0pt}
\setlength{\parskip}{6pt}
\setlist{itemsep=3pt,topsep=5pt,leftmargin=*}
\setlength{\emergencystretch}{2em}
\DefineVerbatimEnvironment{LeanCode}{Verbatim}{fontsize=\small,breaklines=true,samepage=true,
  breakanywhere=true,frame=single,rulecolor=\color{accent},framesep=5pt}
\newcommand{\Commit}{60e9c5d620c6c387d932876cfa2d9376f64ea56c}
\newcommand{\src}[2]{\href{https://github.com/VladimirReshetnikov/ProveIt/blob/60e9c5d620c6c387d932876cfa2d9376f64ea56c/Computability/DiophantineConsistency/Lean/DiophantineConsistency/#1.lean\#L#2}{\texttt{\detokenize{#1}}:L#2}}
\newcommand{\artifact}[1]{\nolinkurl{#1}}
\newcommand{\status}[1]{\textbf{#1}}

\begin{document}
\begin{titlepage}
\vspace*{0.7in}
{\Large\sffamily\color{accent} SOURCE REVIEW AND CHECKED REFACTORINGS\par}
\vspace{0.35in}
{\Huge\sffamily\bfseries\color{ink} Shortening and simplifying\par
the DiophantineConsistency\par Lean proofs\par}
\vspace{0.45in}
{\large Prepared for Vladimir Reshetnikov\par}
{\large 26 October 2026\par}
\vspace{0.4in}
\begin{tabularx}{\textwidth}{@{}lX@{}}
\toprule
Repository & \url{https://github.com/VladimirReshetnikov/ProveIt}\\
Reviewed commit & \texttt{60e9c5d620c6c387d932876cfa2d9376f64ea56c}\\
Lean & \texttt{leanprover/lean4:v4.32.0}\\
Mathlib & \texttt{81a5d257c8e410db227a6665ed08f64fea08e997}\\
Scope & \texttt{Computability/DiophantineConsistency}\\
\bottomrule
\end{tabularx}
\vfill
\colorbox{pale}{\parbox{0.94\textwidth}{\textbf{Main recommendation.}
Concentrate on repeated proof structure and missing shared interfaces:
uniform coefficient bounds, library lookup facts, successful decoder cases,
assignment transport, finite residual sums, and scanner invariants.
Preserve the executable expressions, public statements, and each theorem's
axiom dependencies while reducing the source burden.}}
\vspace{0.2in}
\end{titlepage}

\tableofcontents
\newpage

\section{Results and priorities}

The strongest opportunities are concrete and local. Several long passages
repeat the same mathematical argument with different indices or repeat
facts already available in Lean, Mathlib, or a neighboring ProveIt module.
The proposed changes therefore start by making those facts reusable. A
wholesale change of the mathematical encoding would add considerably more
verification work than these proof refactorings require.

This package separates three kinds of evidence. A \emph{full-file replay}
checks a replacement file with its original imports and declarations. A
\emph{checked prototype} establishes a proposed reusable lemma or a copied
original theorem statement, without claiming that every client module has
been migrated. An \emph{architectural proposal} identifies a source-specific
change that still needs implementation and measurement.

\begin{longtable}{@{}>{\raggedright\arraybackslash}p{0.24\textwidth}>{\raggedright\arraybackslash}p{0.38\textwidth}>{\raggedright\arraybackslash}p{0.31\textwidth}@{}}
\toprule
Change & Concrete benefit & Evidence in this package\\
\midrule
\endhead
Uniform row bounds & Replace 22 repeated coefficient proofs and their dispatch
with one finite-case proof. & Full replacement file: 679 to 122 lines; slower in the recorded elaboration run.\\
Core list lookup lemmas & Seven unchanged declarations across five files:
61 to 21 physical lines. & Core wrappers and a full ProofChecker replay.\\
Decoder success cases & Replace nested \texttt{none}/\texttt{some} trees with
two uses of the successful-bind characterization. & Full replacement file; unchanged per-theorem axiom lists.\\
Assignment transport & Three representation constructors share two lemmas;
260 to 242 lines, including the helpers. & Full replacement file and
unchanged-expression checks.\\
Finite residual API & Three compilers share zero-sum and degree arguments;
two also share witness regrouping. & Eight exact replacements; four patched modules and their selected audits replayed.\\
Scanner and history APIs & Share recursive-step, terminal-state, append,
and invariant arguments. & Generic lemmas and specialization bridges checked.\\
Small arithmetic changes & Reuse finite-sum, division/remainder, power, and
cast results with the same hypotheses. & Checked exact-statement examples.\\
Object semantics and generated circuits & Reduce completeness boilerplate,
reuse degree interfaces, and certify power tables. & Source-specific proposals;
integration remains future work.\\
\bottomrule
\end{longtable}

The source-line measurements are deliberately narrow. They include helper
definitions when a helper is part of a proposed patch. They are not estimates
of a repository-wide percentage reduction, proof-term size, compilation
speed, or the size of the resulting Diophantine equation. The equation data
remain unchanged in the accepted local refactorings.

\subsection{What the survey covers}

The pinned subtree contains 711 tracked Lean files. The maintained
\texttt{Lean} tree contains 546 files; 363 of them lie directly in the
\texttt{Lean/DiophantineConsistency} module directory, totaling 68,308 physical
lines. That directory includes 38 files with a \texttt{Checks} or
\texttt{Audit} suffix. Counting the separate \texttt{Lean/checks} directory
gives 220 audit/check files in the maintained tree. The remaining tracked
Lean files in the requested subtree are in historical proposal/provenance
and candidate directories. They should not be counted as additional live
implementations to refactor automatically.

The complete filename, byte count, line count, and SHA-256 inventory is in
\artifact{evidence/source-inventory.csv} and its JSON counterpart. These
numbers come from the supplied \artifact{scripts/inventory.py}. Its tactic
and duplicate-name counts are textual indicators; they are not Lean AST
analysis. In the main module directory, after removing comments and strings,
the survey found no \texttt{sorry}, \texttt{admit}, \texttt{axiom}, or
\texttt{native\_decide} tokens. This is not a claim that imported theorems
have no axioms. The separate kernel dependency checks matter precisely
because a textual scan cannot establish that.

\section{Replace repeated coefficient bounds with one argument}
\label{sec:coefficients}

\subsection{The existing pattern}

\src{NativeCoefficientBounds}{63} begins a sequence of 22 private
\texttt{rowEnvelope\_i} lemmas. Each lemma fixes one row, unfolds its literal
coefficient recipe, and proves that every branch of a nested conditional is
bounded by the same envelope. The public dispatcher starts at
\src{NativeCoefficientBounds}{602}. The repeated block contains 566 physical
lines, including spacing and its explanatory text.

For ports $p_0,p_1,p_2,p_3\geq0$, the common envelope is
\[
 E=2p_0^2+2p_3p_0+p_2+p_3+2p_0+p_1p_0+p_1+8.
\]
Every branch proof uses exactly the same facts: the four nonnegativity
hypotheses, $p_0^2\geq0$, $p_3p_0\geq0$, and $p_1p_0\geq0$. Once these
products are treated as known nonnegative terms, the final obligations are
linear arithmetic. There is no need to copy the entire row expression into
a separate \texttt{heq : ... := rfl} in every lemma.

\subsection{The replacement}

The important operation is to apply the existing conditional-bound lemma
uniformly, rather than split on every equality test and accumulate unused
branch hypotheses. The core of the candidate is:

\begin{LeanCode}
have hsq := sq_nonneg (ports 0)
have hmf := mul_nonneg h3 h0
have hkp := mul_nonneg h1 h0
fin_cases k <;> dsimp only [rowCoefficient]
all_goals
  unfold envelope
  repeat' apply abs_ite_le
  all_goals
    rw [abs_le]
    constructor <;> linarith only [h0, h1, h2, h3, hsq, hmf, hkp]
\end{LeanCode}

Here \texttt{abs\_ite\_le} is the existing private helper: if both branch
values have absolute value at most $E$, then so does the conditional. The
proof is mathematically identical to the original strategy. The shorter
script states that strategy once and lets it act on all 22 literal rows.

The full replacement file compiles against the original imports. Excluding the review-only axiom-print commands, it shrinks from 679 to 122 physical lines, saving 557 lines (about 82\%). The repeated row block shrinks from 566 to 15 lines, including a local \texttt{maxHeartbeats} option; a separate symbolic-power simplification accounts for the remainder. The three audited exported theorems retain their original standard axiom lists. The coefficient table, envelope definition, and all public theorem signatures are unchanged. The patch is supplied as an optional source-maintenance tradeoff in \artifact{patches/optional/coefficient-bounds-compact.patch}.

The public assumptions, including all four nonnegativity hypotheses, are
essential to this argument and remain present. The patch does not assert
that the bound holds for arbitrary signed ports, and it does not alter the
coefficient table or the stated published height.

\subsection{Proof budget and maintenance}

Combining these small proofs into one does change the heartbeat budget: the default 200,000 limit timed out. The checked candidate uses a local \texttt{set\_option maxHeartbeats 4000000 in} on the combined theorem. This makes its proof budget explicit. It should not become
an unexplained file-wide or project-wide increase. A proposed alternative retaining separate row declarations was also tested, but hit the default heartbeat/simplifier limits on a large row. It is not included among the checked replacement files. A future shared branch-proof recipe or checked coefficient certificate is preferable to pretending that this fallback is already validated.

Source reduction alone does not show smaller proof terms or faster elaboration. In one sequential comparison using one Lean thread, the original file used 23.37 seconds of user CPU time and the compact file used 110.40 seconds; wall times were 24.42 and 111.37 seconds. Both runs exited successfully under the same 3,000 MB Lean memory cap. These are single-run diagnostics on a shared host, not a general benchmark, but the observed roughly 4.7-fold CPU increase is enough to reject an unqualified performance improvement claim. Prioritize the other checked patches when build time matters; use this compact version when the reduced proof-maintenance burden is worth that tradeoff.

\subsection{A certificate design that could remove repeated arithmetic search}

The common envelope is a nonnegative linear combination of a fixed eight-entry basis:
\[
 v=(p_0^2,\ p_3p_0,\ p_2,\ p_3,\ p_0,\ p_1p_0,\ p_1,\ 1),
 \qquad b=(2,2,1,1,2,1,1,8).
\]
For each literal branch, record an integer coefficient vector $q$ with
$|q_i|\leq b_i$. A single triangle-inequality theorem gives
\[
 \left|\sum_i q_i v_i\right|\leq\sum_i |q_i|v_i
 \leq\sum_i b_i v_i=E.
\]
A generator could emit the small vectors, ordinary decidable componentwise
bound certificates, and kernel-checked equalities connecting their
evaluations to the actual row expressions. This would remove repeated
linear-arithmetic search, rather than merely its repeated source spelling.
The exact binding to the existing coefficient definition remains an
obligation. This certificate architecture is a proposal, not an implemented
or measured improvement in this package.

\section{Use existing lookup, sum, and arithmetic lemmas}

\subsection{Seven list inductions have direct replacements}

Five modules independently prove that a successful optional list lookup
returns a member of the list. Two also prove the converse existence of an
index for a member. The exact core lemmas are
\texttt{List.mem\_of\_getElem?} and \texttt{List.getElem?\_of\_mem}, from the
pinned Lean release's \texttt{Init/Data/List/Lemmas.lean}.

\begin{LeanCode}
theorem lookup_mem {α : Type} {xs : List α} {n : Nat} {x : α}
    (h : xs[n]? = some x) : x ∈ xs :=
  List.mem_of_getElem? h

theorem exists_lookup {α : Type} {xs : List α} {x : α}
    (h : x ∈ xs) : ∃ n, xs[n]? = some x :=
  List.getElem?_of_mem h
\end{LeanCode}

The patch preserves the original local declaration names and theorem
statements. It changes only their proof bodies. Its seven complete
declarations go from 61 to 21 physical lines. The sites are
\src{AnnotatedAssumption}{40}, \src{AnnotatedRuleSide}{18},
\src{AnnotatedBound}{50}, \src{AnnotatedRename}{110}, and
\src{AnnotatedProofChecker}{177}. The last file was replayed in full, with
axiom checks on the changed lemma and on its soundness/checker endpoints.
The other four clients are supported by the exact core-lemma prototypes
and a clean patch application check; their entire dependency closures have
not been rebuilt for this package.

This is a good first integration change: it removes bespoke inductions
without creating a new project helper, import layer, or computational
definition.

\subsection{A finite sum at most one}

\src{CurrentCellLabels2011}{131} proves that a natural-valued finite function
whose sum is at most one is either identically zero or has one entry equal
to one and all other entries zero. Mathlib already characterizes this
situation with \texttt{Finset.sum\_le\_one\_iff}. The nearby module
\src{NativeSelectorClause54}{65} already uses that lemma.

The checked replacement reduces the original 27-line proof body to 10
lines. It uses the library characterization, splits the zero case from a
nonzero entry, and reconstructs the exact original conclusion. The type
remains an arbitrary \texttt{Fintype}; the argument also handles an empty
type. This is a particularly useful style of simplification: transfer an
established local method to another development, instead of adding a
parallel lemma with a slightly different statement.

\subsection{Natural subtraction and remainder}

In \src{CurrentCellTranslation2011}{17}, the whole-cell portion of a shift is
\texttt{t - t \% 2011}. Two short facts can be stated directly using core
division/remainder identities. In particular,

\begin{LeanCode}
theorem wholeShift_add_remainder (t : Nat) :
    (t - t % 2011) + t % 2011 = t :=
  Nat.sub_add_cancel (Nat.mod_le t 2011)
\end{LeanCode}

The inequality is still explicit in the proof term. The argument does not
treat truncated natural subtraction as subtraction in a group. The
companion divisibility example uses
\nolinkurl{Nat.mul_div_self_eq_mod_sub_self} to supply the quotient
witness directly. Both examples compile with only the standard library.

\subsection{Let \texttt{omega} handle the cast interface it already supports}

\src{NativePellStepDown}{161} contains \texttt{eq\_of\_residues}, whose final
reasoning is Presburger arithmetic relating natural bounds and integer
casts. The checked small example removes four manually introduced cast
facts and lets \texttt{omega} use the original hypotheses directly. The complete original and replacement files both compile with matching audit output for the changed theorem and two source-index endpoints. The replacement retains every statement and reduces the file from 248 to 241 physical lines, excluding appended audits. An initially imposed 3,000 MB Lean cap stopped both files during import; both passed when replayed sequentially with a 6,000 MB cap. This was an environment limit shared by the original and replacement, not a failed arithmetic argument.

This is a targeted use of automation. It is not a recommendation to replace
nonlinear Pell arguments with an unstructured \texttt{omega} or
\texttt{nlinarith} call. It works at a specific interface where the
arithmetic is already linear and the cast relations are supported.

\subsection{Keep huge powers symbolic}

The published-scale lemmas in \src{NativeCoefficientBounds}{647} already
avoid expanding million-bit integers. That is the right design. The
\texttt{published\_large\_scale} proof can be shortened further by expressing
$16$ as $2^4$, combining powers with \texttt{pow\_add}, and invoking
\texttt{pow\_le\_pow\_right₀} for the exponent comparison. This preserves the
symbolic proof strategy; it does not use a giant numerical normalization.

\section{Shorten decoder proofs without enlarging their axiom dependencies}

\subsection{Repeated successful-decode cases}

\src{AnnotatedProofCode}{78} and \src{AnnotatedProofCode}{107} prove two
properties of a successful fuel-bounded tree decoder:
\[
 \operatorname{decode}(f,c)=\operatorname{some}(t)
 \quad\Longrightarrow\quad
 \operatorname{depth}(t)\leq f,
 \qquad \operatorname{quote}(t)=c.
\]
The inductions are over the same fuel parameter. In the fork case, both
proofs explicitly split the left recursive decode into \texttt{none} and
\texttt{some}, then do the same for the right decode. The unsuccessful
branches contradict the successful outer result.

The useful abstraction is the successful-bind characterization for
\texttt{Option}. After reducing the decoder branch, apply
\texttt{Option.bind\_eq\_some\_iff.mp} directly to the hypothesis. Obtain the
left decoded tree and its equation, then obtain the right decoded tree and
its equation. The final \texttt{some} equality identifies the output tree.
This removes the impossible branches from the source proof.

Direct application is significant here: a \texttt{simp} rewrite with the
same theorem did not match one already-reduced bind shape in the tested
development. The supplied file records the working elaborated version.

Both the isolated original decoder definitions and the complete modified \texttt{AnnotatedProofCode} file compile. The baseline and replacement audit output is byte-for-byte identical for the two directly changed theorems and both checked checker endpoints. The recommended patch is \artifact{patches/decoder-options.patch}.

\subsection{Why the apparently shorter fusion is not the default patch}

A combined theorem asserting both conclusions shortens the original
55-line block to 35 lines, with the old names as projections. That version
does compile. However, the original depth-bound theorem depends on
\texttt{propext} and \texttt{Quot.sound}; the quotation theorem additionally
depends on \texttt{Classical.choice}, inherited through
\texttt{Nat.pair\_unpair}. Projecting the depth bound from a proof of the
conjunction makes the depth theorem inherit that extra axiom too.

The accepted approach keeps the inductions independent and compresses
their \texttt{Option} reasoning. Its block is 43 lines, saving 12, while
retaining the original dependency lists:

\begin{center}
\begin{tabular}{@{}lp{0.63\textwidth}@{}}
\toprule
Theorem & Original and accepted replacement axioms\\
\midrule
\texttt{decoded\_depth\_le} & \texttt{propext}, \texttt{Quot.sound}\\
\texttt{quote\_of\_decode} & \texttt{propext}, \texttt{Classical.choice}, \texttt{Quot.sound}\\
\bottomrule
\end{tabular}
\end{center}

The final checker endpoints already used the standard classical axioms,
so comparing only those endpoints would have missed the local change.
This example gives a practical reason to audit the directly refactored
theorems as well as their main consumers. The fused prototype is included
as an explicitly non-default alternative, with its tradeoff recorded.

\section{A small assignment-transport API removes repeated extensionality}

\subsection{The invariant already implicit in the constructors}

\src{Representation}{155} defines \texttt{Rep.reindex}, followed by
conjunction and disjunction constructors. They preserve executable syntax
and package correctness proofs with an explicit number of witnesses.
Several local proofs establish the same identity for the two components
of a sum-indexed assignment.

For input assignments $v$, witness assignments $w$, and renamings $f,g$,
the common fact is
\[
 \operatorname{natAssignment}(v,w)\circ\operatorname{Sum.map}(f,g)
 =\operatorname{natAssignment}(v\circ f,w\circ g).
\]
Its proof is function extensionality followed by the two constructors of
\texttt{Sum}. Combining it with the existing \texttt{Expr.eval\_rename}
gives the evaluator-level interface:

\begin{LeanCode}
theorem eval_rename_natAssignment (e : Expr (σ ⊕ α))
    (f : σ → τ) (g : α → β) (v : τ → Nat) (w : β → Nat) :
    (e.rename (Sum.map f g)).eval (natAssignment v w) =
      e.eval (natAssignment (v ∘ f) (w ∘ g)) := by
  rw [Expr.eval_rename, natAssignment_map]
\end{LeanCode}

The correctness field of reindexing then becomes:

\begin{LeanCode}
correct v := by
  simpa only [eval_rename_natAssignment, Function.comp_id]
    using r.correct (v ∘ f)
\end{LeanCode}

Conjunction and disjunction use the same lemma for the left and right
blocks of a finite witness vector. Their remaining proof is the actual
logical work: combine or separate the witnesses, then use the zero-sum-of-
squares or zero-product characterization.

\subsection{Data preservation is checked, not inferred from equivalence}

The complete replacement file compiles. Its size changes from 260 to 242
physical lines, including 15 new helper lines. The separate prototype also
checks by \texttt{rfl} that all three new constructors have exactly the same
\texttt{expression} fields as the original constructors. Their witness
counts are unchanged by construction.

This distinction matters in this repository. A shorter proof of a merely
equivalent solvability proposition would not by itself justify claims
about the same executable equation or the same witness count. Here the
data fields are preserved, and only the supporting correctness proofs are
reorganized. Keep these helpers near \texttt{natAssignment} in the
representation layer so that later compilers can reuse them without
introducing reverse imports.

\section{Share the finite residual and degree arguments}

\subsection{Three copies of the same evaluator lemma}

The private \texttt{eval\_sumExpressions} induction appears in
\src{DAGQuartic}{83}, \src{NaturalDAGQuartic}{110}, and
\src{QuadraticRegisterSystem}{70}. Each client then proves the same
zero-sum-of-squares characterization for its residual expressions.

The shared interface should take arbitrary finite residual syntax
$e_i:\operatorname{Expr}(\sigma)$ and an integer assignment $v$:
\[
 \operatorname{eval}\!\left(\sum_{i:\mathrm{Fin}\,n}e_i^2,v\right)=0
 \quad\Longleftrightarrow\quad
 \forall i,\ \operatorname{eval}(e_i,v)=0.
\]
Make the evaluation lemma reusable, then prove this characterization once.
Each of the three client theorems becomes an instantiation. The $n=0$ case
is covered by the finite statement: the empty sum is zero and the universal
condition is vacuous.

Keep the domain assumptions honest. The proposed lemma is about integer
evaluation. The corresponding assertion is false in an arbitrary ring;
for example, nonzero complex squares can cancel. A more general theorem
would need explicit ordered-ring or suitable sum-of-squares hypotheses.

\subsection{The needed polynomial-degree theorem already exists}

For residual degrees at most $d$, the compiled sum of squares has total
degree at most $2d$. The H10 common library already supplies
\texttt{Diophantine.totalDegree\_sum\_pow\_le}. The proposed shared helper
uses it rather than proving another polynomial-degree induction.
The three degree clients are \src{DAGQuartic}{244},
\src{NaturalDAGQuartic}{253}, and \src{QuadraticRegisterSystem}{150}.

This is an upper bound, not an assertion that the degree is exactly $2d$.
Cancellations or zero residuals can lower the degree. The proposed API
retains the original inequalities.

\subsection{Regroup witnesses with the existing sum-type theorem}

The proofs at \src{DAGQuartic}{260} and \src{NaturalDAGQuartic}{267}
manually split an assignment on a sum type into its two components,
reconstruct it with \texttt{Sum.elim}, and prove cast identities by
extensionality. Mathlib's \texttt{Sum.exists\_sum} already gives the
existential regrouping. A single shared cast identity then allows both
clients to use their existing \texttt{exists\_zero\_iff} theorem through
\texttt{exists\_congr}.

The topological, allowed-node, and output-index hypotheses are retained.
The shorter proof does not assume the evaluation/constraint correspondence;
it invokes the existing theorem that establishes it.

The checked patch adds a 36-line \texttt{CompilerSupport} module and changes the three client files from 285 to 257, 293 to 264, and 170 to 158 physical lines. The net saving, including the new module, is 33 physical lines (36 nonblank lines). All eight original/replacement theorem pairs have the same audited standard axioms. An isolated overlay also compiled all four patched modules, eight unchanged dependencies needed by the audits, and all three existing audit files: 15 successful commands. The overlay does not modify the original checkout. It is a targeted integration replay, not a complete downstream rebuild.

\subsection{Two further consolidations}

\textbf{Generic sum cost.} The private
\texttt{DAGQuartic.sumExpressions\_cost} is general in a per-expression
bound $b$. \texttt{QuarticCost.sumExpressions\_binaryOps\_le} separately
proves the case $b=5$. Move the generic statement into the common cost
module, keep the old specialized name as a compatibility wrapper, and
instantiate the signed DAG's bound at $b=11$. Directly substituting the
existing $b=5$ theorem for a $b=11$ use would be incorrect.

\textbf{Weighted syntax degree.} \texttt{QuadraticRegisterSystem.syntaxDegree}
is the constant-weight-one case of the existing
\texttt{Expr.weightedDegree} in \src{WeightedExpressionDegree}{1}.
Its polynomial-degree proof can be derived from
\texttt{degree\_evalRing\_le} and \texttt{evalRing\_X}, after proving an
API-preserving equality for the syntax measure. Check the import graph
before moving this fact; a small common degree module is preferable to a
cycle between compiler clients.

These last two changes are proposals in this package. They should be
implemented after the smaller checked interfaces, not combined with an
unrelated redesign of the expression datatype.

\section{Reuse scanner and validated-history proofs}

\subsection{Recursive steps and terminal states}

The annotated scanner modules repeat proofs about a \texttt{Nat.rec}
iteration: its zero case, its successor equation, and the fact that a
finished or rejected state stays fixed. Relevant families include
\texttt{AnnotatedBoundScan}, \texttt{AnnotatedLineScan},
\texttt{AnnotatedRenameScan}, \texttt{AnnotatedAxiomContext}, and
\texttt{AnnotatedContext}. Representative anchors are \src{AnnotatedBoundScan}{30}, \src{AnnotatedLineScan}{35}, \src{AnnotatedRenameScan}{34}, \src{AnnotatedAxiomContext}{32}, and \src{AnnotatedContext}{100}.

The checked \artifact{lean/ScanIteration.lean} prototype proves generic
recursive-step and fixed-state facts and checks concrete scanner shapes.
One can reuse these proof lemmas while leaving every executable scanner
definition intact. That is a useful first step because scanner arithmetic,
primitive-recursive encodings, and size theorems may depend on the literal
definition shape.

\subsection{An input row and a stored judgment need not have the same type}

\texttt{AnnotatedProofChecker}, \texttt{AnnotatedRename}, and
\texttt{AnnotatedBound} repeat an append theorem, soundness induction, and
\texttt{Produces.append}/\texttt{Produces.step} arguments for a validator
that pushes a new item onto a history. The reusable runner takes:

\begin{itemize}
\item an input type $I$ and a stored-history type $S$;
\item a validator $I\to\operatorname{List}(S)\to\operatorname{Bool}$;
\item an output map $I\to S$;
\item a local soundness theorem for a predicate on stored items.
\end{itemize}

The generic invariant states that a successful run preserves that predicate
throughout the history. The output map is necessary: a proof-checker line
may store a judgment rather than the entire input line. A shared runner
with a single type would not cover all three clients accurately.

The prototype \artifact{lean/ValidatedHistory.lean} checks the generic
runner, all four reusable theorems, and bridges from unchanged concrete
recursion shapes. Those bridges are proved function equalities; the
concrete recursive definitions are not automatically definitionally equal
to a specialization of a separately defined generic runner. Replacing
them by an unjustified \texttt{rfl} would fail.

For integration, first keep the concrete definitions and use the proved
bridges in theorem wrappers. Only consider redefining the executable
runner itself after checking the downstream numeric and primitive-recursive
representations. The generic prototypes are checked; migration of every
concrete client remains future work.

\section{Simplify object-theory endpoints at the semantic boundary}

\subsection{A completeness-and-sealing helper}

Many \texttt{Object*} endpoints follow the same sequence: invoke first-order
completeness, introduce an arbitrary model and its ZFC-model hypothesis,
remove the empty assumption context, rewrite the sealed formula's
semantics, and apply a previously established semantic lemma. The survey
found 27 completeness invocations in these files, with two notably
different shapes. About 25 are plausible clients for a helper that turns
validity of the unsealed body over arbitrary models into a ZFC derivation of its seal. One representative endpoint is \src{ObjectNaturalExpression}{182}.

The candidate helper belongs in a low common bridge layer. It must quantify
over the same arbitrary models, environments, and formula semantics as
the existing completeness theorem. It should encapsulate only the repeated
completeness/sealing steps. The arithmetic content remains the semantic
theorem supplied by the caller.

This is not permission to replace the object model's natural-number domain
by the metatheory's \texttt{Nat}. The development must still handle its full
internal $\omega$, including nonstandard behavior allowed by the model
semantics. The existing generic semiring interface in
\texttt{ObjectNaturalSemiring.evalRel\_iff} is the appropriate kind of
abstraction downstream. Using it in its own foundational dependency
\texttt{ObjectNaturalExpression} would create an import cycle in the current organization unless the foundational definitions were split deliberately.

\subsection{Reuse the ZFC-model-to-ZF-axioms bridge}

Two copies of the six-axiom extraction are \src{ObjectArithmetic}{143}
and \src{ObjectMultiplication}{169}.
There is also an existing theorem
\nolinkurl{LeanProofs.BoundedZFCConsistency.zfAxioms_of_zfcModel} in
\texttt{CodedRank}. Extract a suitable common theorem or expose a low-layer
alias. Importing the large later \texttt{CodedRank} development into an
earlier arithmetic foundation simply to save these lines would be a poor
dependency tradeoff.

\subsection{A smaller quantified-list example}

\texttt{ObjectNaturalExpression.}\allowbreak\texttt{inputsNaturalForm\_spec} manually constructs
and destructs membership in \texttt{List.ofFn}. Mathlib already has
\texttt{List.forall\_mem\_ofFn\_iff}. The anticipated proof is a short
\texttt{simp only} chain through the conjunction semantics, that list
theorem, and \texttt{fNatF\_spec}. This is a precise next candidate, but its
complete original-import replay is not included here; treat it as a
proposal rather than a checked patch.

\section{Unify modular and finite-word interfaces}

\subsection{Share residue transport at the correct modulus}

\src{CurrentClause2011}{27} has public addition/subtraction residue
theorems whose proof bodies are duplicated by private helpers in
\src{CurrentReflectedGeometry2011}{74}. Put the common facts in a lightweight
cyclic-residue module imported by both clients. A checked generalized
prototype appears in \artifact{lean/SmallArithmetic.lean}.

The modulus must remain the full cyclic modulus used in the application,
such as $N L$. Replacing it everywhere by the local cell length $L$ loses
whole-cell position information. A small lemma transporting equality of
residues is safe; silently quotienting away a whole-cell offset is a change
of theorem.

\subsection{One boundary between \texttt{ZMod}, \texttt{Fin}, and lists}

\src{NativeCyclicConvolution}{1} and \src{NativeRadixDigits}{1} repeatedly
move between cyclic functions, finite functions, padded digit lists, and
finite sums. A common API should contain the evaluation of a word built
from \texttt{List.ofFn}, the lookup identity through the \texttt{Fin}/
\texttt{ZMod} equivalence, and the finite-sum transport theorem. Then word
reconstruction and mask statements can be proved at the same abstraction
level before crossing that boundary once.

The existing \texttt{word\_eq\_ofDigits}, \texttt{ofDigits\_ofFn}, and
\texttt{finEquiv\_val} lemmas already do part of this work. The proposal is
to provide the missing composition as a shared interface, not introduce a
second digit representation. Preserve the \texttt{NeZero} modulus instance,
digit-capacity bounds, and padded-list length hypotheses at the boundary.

\subsection{Package repeated geometric facts without adding assumptions}

The current-cell and reflected-geometry proofs repeatedly derive positivity
of a modulus, source-index bounds, and wrapped/nonwrapped index equations.
A small record of \emph{proved consequences} can shorten subsequent
calculations. Its constructor should derive those fields from the existing
parameters; it must not turn an unproved global alignment or label-binding
fact into an implicit typeclass assumption. Keep these records explicit
arguments when two different geometric situations may coexist.

This remains a design proposal. The immediate residue and natural-remainder
examples are much lower-cost starting points.

\section{Treat generated circuit proofs as generated interfaces}

\subsection{What the large files actually contain}

The two generated step files dominate the raw line-count list:

\begin{center}
\begin{tabular}{@{}lrrr@{}}
\toprule
File & Physical lines & Step lemmas & Power lemmas\\
\midrule
\texttt{GapCircuitSteps66} & 3,122 & 484 & 86\\
\texttt{SplitCircuitSteps59} & 3,445 & 515 & 107\\
\midrule
Total & 6,567 & 999 & 193\\
\bottomrule
\end{tabular}
\end{center}

Each step lemma already instantiates the shared
\texttt{run\_node\_equation}. These are not 999 distinct missing mathematical
arguments. A source-volume improvement would change the generator and its
public interface: expose one indexed theorem, then keep numerical aliases
only where clients or audit scripts need stable names.

Do not hand-edit generated output and leave the generator unchanged. The
report's proposal targets the step generators as well as their output,
with a regeneration check in the integration gate.

\subsection{Reuse the existing power-certificate theorem}

The 193 power lemmas contain 193 local \texttt{ring} invocations. The
existing \texttt{ArithmeticDAG.run\_power} in \src{ArithmeticDAGPowers}{24} already proves correctness
of exponent labels under \texttt{PowerStep} conditions. The generators
already compute the relevant exponent information. A useful next design is:

\begin{enumerate}
\item generate one exponent-label table per circuit;
\item certify its local \texttt{PowerStep} obligations in Lean;
\item invoke \texttt{run\_power} once;
\item recover any required \texttt{natural\_power\_*} names as short aliases.
\end{enumerate}

This is a potential reduction in repeated normalization, but it is not
implemented or benchmarked in this package. The size and elaboration cost
of the checked table must be measured against the original 193 proofs.
Use ordinary kernel-checked certificates; substituting
\texttt{native\_decide} would change the verification mechanism and is not
an automatic source simplification.

\subsection{A naming trap worth avoiding}

The names \texttt{naturalNodes} in the Gap/Split developments describe
natural inputs. They do not imply that every intermediate node satisfies
\texttt{NaturalDAGQuartic.Allowed}. The Gap list contains subtraction nodes
and negative literals. A shared natural-semiring proof must not be applied
solely because the identifier contains ``natural.'' Retain the explicit
allowed-node certificate wherever that theorem requires it.

\section{A reusable invariant for arithmetic DAG evaluation}

Two existing proofs use nearly the same strong induction along earlier
dependencies: \nolinkurl{NaturalDAGQuartic.run_nonnegative} and
\nolinkurl{ArithmeticDAG.value_degree_le}. A generic
\texttt{run\_invariant} can take a predicate $P(i,v)$ and a local closure
rule for each node given the invariant on its dependencies. Nonnegativity
and degree bounds then become local closure proofs.

The repository already has \texttt{run\_node\_equation} and
\texttt{run\_eq\_of\_node\_equations}; those theorems are the natural basis
for such an interface. Do not duplicate their semantic correspondence in
a new evaluator. A small common induction lemma is sufficient.

The \artifact{lean/RunInvariant.lean} prototype compiles, including the shared induction rule and both exact application statements. The two original/replacement application pairs have matching axiom lists; all directly printed lists contain only the standard axioms. This establishes the reusable proof interface; moving it to its final import layer and replacing the original two proof bodies remain an integration step.

A more ambitious alternative would make the natural-register evaluator
semiring-polymorphic and relate it to the erased integer DAG by an explicit
erasure-correctness theorem. This could reduce cast and nonnegativity
proofs, but would also introduce new dependent-index transports. It is a
later research direction, not the recommended first refactoring.

\section{Validation and integration plan}

\subsection{What has been checked}

The evidence files record the exact compiler version, exit status, source
hashes where supplied, and axiom output for the tested examples. A successful audit requires compiler exit status zero and no proof errors before its axiom output is considered. In one timed-out trial, an axiom-printing command still emitted an innocuous-looking line; that line did not make the failed declaration a valid proof. All
accepted checked examples use ordinary Lean proofs. The local source
refactorings do not add an admitted theorem or a new axiom. The rejected
default decoder fusion is retained separately so that its verified
\emph{increase in the local axiom dependency list} is visible.

The completed checks include the two syntax replacements with paired full-file axiom logs; the assignment patch and its executable-data equality checks; the finite-residual patch with its 15-command overlay replay; the DAG invariant prototype; the full original/compact coefficient files; the full original/shortened Pell-step files; and the small core/arithmetic prototypes. Across the structural examples, 13 original/replacement theorem pairs have equal audited axiom sets. The \texttt{verification-plan.json} file and topic receipts give the exact scope for each item.

The environment was shared and resource-constrained. Timings are single-run diagnostics rather than reliable general performance estimates. The observed coefficient-proof slowdown is reported explicitly because it affects the integration recommendation; no speedup is claimed. The initial broad prerequisite build was replaced by a
sequential build of the needed missing modules. The compiler inputs remained
the pinned Lean release, the pinned Mathlib revision, and the original
ProveIt sources. Toolchain/cache preparation is not included in the ZIP.

\subsection{Replay the evidence}

\artifact{scripts/verify.py} is the entry point for replaying the packaged
examples against a checkout at the pinned commit. The accompanying
\artifact{README.txt} specifies prerequisites, the available verification
levels, and which artifacts are proposals. The script checks the recorded
source snapshot and compiles the requested files; it does not fetch a
different Mathlib revision or silently accept a newer repository commit.

The report and patches are designed for review and integration. A full
build of all 546 maintained Lean files, or every downstream consumer of
the combined patch, is outside the validation claim of this package.
The next integration gate should replay affected import dependents,
especially where a private helper is removed, a shared lemma is introduced,
or a simplification attribute changes.

\subsection{Suggested order}

\begin{enumerate}
\item Apply the list-lookup and small arithmetic replacements. Keep public
names and statements unchanged and compare the directly changed theorem
axioms.
\item Apply the independent decoder simplification and assignment-transport
patch. Retain the expression-data equality checks for the latter.
\item Review the optional compact coefficient-bound proof separately. Its local heartbeat budget and observed elaboration cost are material tradeoffs, even though it passes the original published-bound statements.
\item Introduce the finite-residual support module and migrate its three
clients together. Count the new module when measuring the reduction.
\item Introduce scanner/history proof interfaces while preserving concrete
definitions through explicit bridges.
\item Review the object-sealing and generated-power designs separately,
because their import/API and regeneration consequences are broader.
\end{enumerate}

For each change, distinguish four properties: the public proposition is
unchanged; executable data are unchanged where relevant; the directly
changed theorem's axiom list has not grown; and affected dependents compile.
One property does not automatically imply the others.

\section{Further questions}

\begin{enumerate}
\item Can a small tool flag repeated extensionality and recursive-case
proof shapes, leaving the choice of abstraction to a maintainer?
\item Can the eight-term coefficient-basis certificate remove the measured elaboration overhead of the compact row-bound proof while retaining its short source and exact literal binding?
\item Does a \texttt{PowerStep} label-table certificate reduce total
elaboration work or serialized proof size relative to the 193 generated
power proofs? Measure both metrics; a shorter generator alone does not
answer the question.
\item How small can the common residual API remain while serving the
signed DAG, natural DAG, and quadratic-register compiler? A good boundary
should avoid importing a full compiler merely to use a finite-sum lemma.
\item Can a successful-decode induction principle reduce both decoder
proofs while preserving their separate constructive/classical dependency
footprints? The conjunction experiment shows exactly what must be avoided.
\item How much object-level boilerplate disappears after one carefully
placed completeness-and-sealing theorem, while preserving all quantification
over arbitrary models and internal $\omega$?
\end{enumerate}

\appendix
\section{Source and artifact guide}

All source links in this report point to the pinned repository revision.
Core-list references are to
\href{https://github.com/leanprover/lean4/blob/8c9756b28d64dab099da31a4c09229a9e6a2ef35/src/Init/Data/List/Lemmas.lean}{Lean's list lemmas at the release commit}.
Mathlib reuse is grounded in the pinned source, including
\href{https://github.com/leanprover-community/mathlib4/blob/81a5d257c8e410db227a6665ed08f64fea08e997/Mathlib/Data/Sum/Basic.lean}{\texttt{Sum.exists\_sum}},
\href{https://github.com/leanprover-community/mathlib4/blob/81a5d257c8e410db227a6665ed08f64fea08e997/Mathlib/Algebra/Order/BigOperators/Group/Finset.lean}{\texttt{Finset.sum\_le\_one\_iff}}, and
\href{https://github.com/leanprover-community/mathlib4/blob/81a5d257c8e410db227a6665ed08f64fea08e997/Mathlib/Data/List/OfFn.lean}{\texttt{List.forall\_mem\_ofFn\_iff}}.

\begin{longtable}{@{}p{0.27\textwidth}p{0.66\textwidth}@{}}
\toprule
Artifact & Purpose\\
\midrule
\endhead
\texttt{*.tex}, \texttt{*.pdf} & This complete report, with editable LaTeX source.\\
\texttt{patches/} & Independent reviewable patches; patch ordering and status
are described in the README.\\
\texttt{lean/} & Checked prototypes and full replacement files, with clear
distinctions between accepted candidates and optional experiments.\\
\texttt{evidence/} & Source inventory, patch metrics, validation receipts,
axiom logs, and a precise status manifest.\\
\texttt{scripts/inventory.py} & Reproduce the maintained-source inventory.\\
\texttt{scripts/verify.py} & Replay selected examples and replacement-file checks.\\
\texttt{README.txt} & Contents, prerequisites, replay instructions,
integration cautions, and status definitions.\\
\texttt{SHA256SUMS.txt} & Integrity hashes of the packaged files.\\
\bottomrule
\end{longtable}

\end{document}
